\documentclass[10pt,a4paper]{amsart}
\usepackage[margin=19mm]{geometry}
\usepackage{amsmath,amssymb,mathtools}
\usepackage[scr=rsfs,cal=euler]{mathalfa}
\usepackage{xfrac}
\usepackage[unicode,hidelinks,bookmarks=false]{hyperref}
\newcommand{\ie}{i.e.\ }
\newcommand{\cf}{cf.\ }
\newcommand{\resp}{respectively\ }
\newcommand{\Subsection}{\subsection}
\providecommand{\nobreakdash}{\nobreak\hbox{-}}
\setlength{\emergencystretch}{2em}
\multlinegap0pt
\usepackage{bm}
\usepackage{bbm}
\usepackage{dsfont}
\usepackage{xspace}
\usepackage{cancel}
\usepackage{mathtools}
\usepackage[shortlabels]{enumitem}
\usepackage{amsthm}
\makeatletter
\@ifundefined{lemma}{\newtheorem{lemma}{Lemma}}{}
\makeatother
\newcommand*{\vecv}{\mathrm{v}}
\newcommand*{\vecw}{\mathrm{w}}
\newcommand{\Pf}{\mathrm{Pfaff}}
\newcommand{\RR}{\mathbb R}
\newcommand{\ZZ}{\mathbb Z}
\title[Distinct Distances on Pfaffian Curves]{Distinct Distances on Pfaffian Curves}
\author{Abhiram Natarajan, Adam Sheffer}
\date{Source: arXiv:2510.04337v1 (CC BY 4.0)}
\begin{document}
\maketitle

\noindent\emph{Source and licence.} Distinct Distances on Pfaffian Curves. Version arXiv:2510.04337v1 (CC BY 4.0), distributed under the Creative Commons Attribution 4.0 International licence, as recorded in the arXiv licence metadata for this exact version and shown on the abstract page https://arxiv.org/abs/2510.04337v1. Full rights notice: licenses/arxiv-2510.04337-CC-BY-4.0.txt.\par\smallskip

\noindent\textit{Editorial scope.} Counted unit: the Lemma whose source label is \texttt{lem:syms} in the article (source0.tex lines 751--754, source label \texttt{lem:syms}), delivered together with the article's own complete original proof of it (source0.tex lines 755--823, one \texttt{proof} environment of 69 lines). No mathematics is written, repaired or re-derived: statement and proof are the article's own text line for line. Required context retained uncounted: le:pfBez (lines 378--383); isofact:classification (lines 730--747); isofact:def (lines 730--747); isofact:group (lines 740--742). Every definition, display and label of those lines is retained unchanged. Omitted from this package and not redistributed: 1--377, 384--729, 743--750, 824--1010. Nothing inside a retained excerpt is omitted or altered: every retained body is delivered exactly as the article's lines are written, so no omission receipt of any kind accompanies this package. Citations: the retained excerpts carry no citation command at all, so no bibliography entry of the article is transcribed and no reference is redistributed. Reference adaptation: none was needed and none was made; every reference inside the delivered unit resolves inside it, so no unresolved reference or citation is produced. Printed slips kept literal: no printed word, symbol, hypothesis or number of the counted statement or of its proof was altered. Layout and class: the article's own class is \texttt{article} with packages including babel, inputenc, amsmath, color, latexsym, amsthm, amssymb, path, enumitem, url, bbm, subcaption, geometry, hyperref; the wrapper is the lane's pinned \texttt{amsart} editorial wrapper at 10pt on A4, which restates the theorem-like environments the retained text needs and adds the article's own macro definitions that the retained text needs (\texttt{\textbackslash Pf}, \texttt{\textbackslash RR}, \texttt{\textbackslash ZZ}, \texttt{\textbackslash vecv}, \texttt{\textbackslash vecw}), with extra wrapper packages (bm, bbm, dsfont, xspace, cancel, mathtools, enumitem). The delivered numbering is the unit's own. Assets: no retained span contains a figure environment, a graphics-inclusion command, a PDF-page inclusion, a table or a TikZ picture; no image file is copied into this package, the package declares no asset dependency and no image file is redistributed with it. Licence: version arXiv:2510.04337v1 of this article is distributed under the Creative Commons Attribution 4.0 International licence, as recorded in the arXiv licence metadata for this exact version and shown on the abstract page https://arxiv.org/abs/2510.04337v1. A full rights notice is delivered at licenses/arxiv-2510.04337-CC-BY-4.0.txt. Characters: every retained excerpt is pure ASCII, so the delivered engine, XeLaTeX with its default Latin Modern text font, reports no missing character.
\begin{lemma} \label{le:pfBez}
Let $Q$ be a Pfaffian chain of order $r$ and chain-degree $\alpha$. 
Consider nonzero $f_1,f_2 \in \Pf_2(\beta,Q)$ of format $(\alpha,\beta,r)$. 
Then, either $Z(f_1)$ and $Z(f_2)$ have a common one-dimensional component or 
$|Z(f_1)\cap Z(f_2)| = O_{\alpha,\beta,r}(1)$.
\end{lemma}

\begin{enumerate}[label=(\textbf{Iso-\arabic*}),ref=(\textbf{Iso-\arabic*})]
    \item \label{isofact:def} The isometries of $\RR^2$ are the transformations 
    %
    \[ h_{A, \vecw}(\vecv) = A\vecv + \vecw,\]
    % 
    where $A \in O_\RR(2)$ and $\vecw \in \RR^2$.
    \item \label{isofact:classification} A complete classification of isometries of $\RR^2$:
    \begin{enumerate}
        \item The \emph{identity} is obtained when $A = I_2$ and $\vecw = 0$. 
        \item A nonzero \emph{translation} is obtained when $A = I_2$ and $\vecw \neq 0$. 
        \item A nonzero \emph{rotation} is obtained when $\det(A)=1$ and $A\neq I_2$. These are the matrices $A = \begin{pmatrix}
            \cos(\theta) & -\sin(\theta) \\ \sin(\theta) & \cos(\theta)
        \end{pmatrix}$ with $\theta \neq 0$. The center of the rotation is $p = (I_2 - A)^{-1}\vecw$ and the angle is $\theta$. We may also write $h(\vecv)=A(\vecv - p) + p$.
        \item A \emph{reflection} is obtained when $\det(A)=-1$ and $A\vecw = -\vecw$. In particular, it is the reflection across line $\frac{\vecw}{2} + \mathrm{ker}(A - I_2)$. We may also write $h(\vecv) = A\left(\vecv - \frac{\vecw}{2}\right) + \frac{\vecw}{2}$.
        \item A \emph{glide reflection} is obtained when $\mathrm{det}(A) = -1$ and $A\vecw \neq -\vecw$. 
    \end{enumerate}
    \item \label{isofact:group} The set of symmetries of a curve forms a group under composition.
\end{enumerate}

\begin{lemma}\label{lem:syms}
Let $Q$ be a Pfaffian chain with chain-degree $\alpha$ and order $r$. 
Consider $f \in \Pf_2(\beta, Q)$ such that $\gamma= Z(f)\subset \RR^2$ is an irreducible Pfaffian curve that is neither a line nor a circle. Then $\gamma$ has $O_{\alpha,\beta,r}(1)$ symmetries.
\end{lemma}

\begin{proof}
Assume for contradiction that $\gamma$ has {\bf a translation symmetry} $T$.
We write $T(\vecv) = \vecv+\vecw$, with $\vecw\in \RR^2\setminus \{0\}$. 
For a point $p\in \gamma$, we have that $f(p) = 0$, which in turn implies that $f(p+\vecw)=0$. 
Inductively, this means that $f(p + n\vecw) = 0$ for all $n \in \ZZ$. 
That is, $\gamma$ has an infinite intersection with the line $\{p + t\vecw : t\in \RR \}$.
By Theorem \ref{le:pfBez}, $\gamma$ is this line, contradicting the assumption.  
  
Next, we assume that $\gamma$ has two {\bf rotation symmetries}.
Following \ref{isofact:def}, we denote these rotations as $h_1=h_{A_1, \vecw_1}(\vecv)$ and $h_2=h_{A_2, \vecw_2}(\vecv)$.
We denote the centers of the rotations as $p_1$ and $p_2$, respectively.
By \ref{isofact:classification}, we have that $h_{i}(\vecv) = A_i(\vecv-p_i) + p_i$. 
We note that $h_{i}(\vecv)^{-1} = A_i^{-1}(\vecv-p_i) + p_i$.
We recall that $A_1,A_2,I_2$ are in the group $O_\RR(2)$, which is Abelian. 
Combining the above leads to
  \begin{align*}
 &(h_2^{-1} \circ h_1^{-1} \circ h_2 \circ h_1)(\vecv) \nonumber \\
 &\qquad \qquad = (h_2^{-1} \circ h_1^{-1} \circ h_2)(A_1(\vecv - p_1) + p_1) \nonumber \\
 &\qquad \qquad = (h_2^{-1} \circ h_1^{-1})(A_2A_1(\vecv - p_1) + A_2 p_1 - (A_2 - I_2)p_2) \nonumber \\
 &\qquad \qquad = (h_2^{-1})(A_2(\vecv - p_1) + (A_1^{-1}A_2 - A_1^{-1} + I_2)p_1 - (A_1^{-1}A_2 - A_1^{-1})p_2) \nonumber \\
 &\qquad \qquad = \vecv + (A_1^{-1} - A_2^{-1}A_1^{-1} + A_2^{-1} - I_2) p_1 - (A_1^{-1} - A_2^{-1}A_1^{-1} + A_2^{-1} - I_2) p_2 \nonumber \\
 &\qquad \qquad = \vecv + (I_2 - A_2^{-1})(I_2 - A_1^{-1})(p_2 -p_1). \nonumber
 \end{align*}

By \ref{isofact:group}, the above transformation is a symmetry of $\gamma$.
When $p_1\neq p_2$, this symmetry is a translation. 
By the first paragraph of the current proof, $\gamma$ cannot have a symmetry that is a translation.
Thus, all rotational symmetries of $\gamma$ have the same center $p$.

Fix a point $q\in \gamma\setminus\{p\}$.  The images of $q$ under the rotational symmetries of $\gamma$ are distinct points on a circle $C$ with center $p$.
By Lemma \ref{le:pfBez}, either $|C \cap \gamma|=O_{\alpha,\beta,r}(1)$ or $C=\gamma$.
By the assumption that $\gamma$ is not a circle, we have that $|C \cap \gamma|=O_{\alpha,\beta,r}(1)$.
This in turn implies that $\gamma$ has $O_{\alpha,\beta,r}(1)$ rotational symmetries. 

Assume for contradiction that $\gamma$ has two {\bf reflection symmetries} with parallel axes of symmetry. 
Combining two such reflections leads to a translation. 
By \ref{isofact:group}, this translation is also a symmetry of $\gamma$, which contradicts the above.
Thus, each two reflection symmetries of $\gamma$ have non-parallel axes.

Consider two reflection symmetries of $\gamma$.
By the preceding paragraph, the two respective axes of reflection intersect in a point $p$. 
Following \ref{isofact:classification}, we may write
%
\[ h_1(\vecv) = A_1\left(\vecv - \frac{\vecw_1}{2}\right) + \frac{\vecw_1}{2} = A_1\vecv + u_1.  \]
Similarly, we write $h_2(\vecv) = A_2\vecv + u_2$.
Since $p$ is fixed by $h_1$, we have that $A_1 p + u_1 = p$.
Since $p$ is also fixed by $h_2$, applying $h_2$ to both sides above leads to
%
\begin{equation*} %\label{eq:TwoRefl}
A_2(A_1 p + u_1) + u_2 = p \quad \Rightarrow \quad (I_2-A_2A_1)p = A_2 u_1 + u_2.  
\end{equation*}

Combining the above implies that
%
\[ (h_2 \circ h_1)(\vecv) = A_2A_1\vecv + A_2u_1 + u_2 = A_2A_1\vecv + (I_2 - A_2A_1)p. \]
%
By \ref{isofact:classification}, $\det(A_1)=\det(A_2)=-1$, so $\det(A_2A_1)=1$.
That is, the symmetry $(h_2 \circ h_1)(\vecv)$ is a rotation with center $p$.
By the above argument for rotation symmetries, the lines of all reflection symmetries of $\gamma$ intersect at $p$.
We fix one such reflection symmetry $h$ and note that combining it with each other reflection symmetry of $\gamma$ leads to a distinct rotation around $p$.
The above analysis of rotation symmetries implies that $\gamma$ has $O_{\alpha,\beta,r}(1)$ reflection symmetries. 

By \ref{isofact:classification}, it remains to consider the case of {\bf a glide reflection symmetry} $h$ of $\gamma$.
We write $h(\vecv) = A\vecv + \vecw$, where  $A\vecw \neq \vecw$ and $A^2=I_2$.
By \ref{isofact:group}, $(h \circ h)(\vecv) = \vecv + (I_2 + A)\vecw$ is a symmetry of $\gamma$.
However, this is a translation, which is impossible. 
We thus conclude that $\gamma$ has no glide reflection symmetries.
%
\end{proof}

\end{document}
